\chapter{风险、失效条件与下一次更新触发器}

\section{重点风险}

风险框架的目的不是罗列宏观风险，而是将每个研究状态绑定到下一次可观察的证据。只要最先到来的实际财报或现金证据不支持当前预告，状态立即从事件验证降为观察；不会以提高倍数或复用旧目标抵消反证。

\begin{riskbox}[高影响风险]
\begin{itemize}
  \item 业绩预告可能无法由实际收入、利润率与经营现金流验证；预告不能年化。
  \item 商品价格、价差、锂盐价格、储能/电池利润率和库存成本可能改变外部预测分母。
  \item 单一券商目标、调整后利润、目标方法未披露或目标分歧，不构成一致预期。
  \item 投资收益、联营收益、投资公允价值变动与并表口径调整可能使表观利润不能直接用于常态估值。
  \item 价格已先行时，正确的风险控制是撤销低位标签，不是提高目标价以维持旧叙事。
\end{itemize}
\end{riskbox}

上述风险首先影响的是“是否允许重新打开目标模型”，而不是即时方向判断。例如雅化、宏桥和亿纬都必须先取得实际H1收入/OCF，而不是直接等待年末利润。

\section{下一轮明确触发器}

\begin{tabularx}{\textwidth}{L{1.7cm}L{3.5cm}X}
\toprule
\textbf{标的} & \textbf{先后验证顺序} & \textbf{状态变化规则} \\
\midrule
002497 & 先等实际H1收入、OCF与预告利润的对应关系；再在Q3/H2检查FY2026利润桥、锂盐量价、库存和套保解释。 & 监测规则：H1 OCF若仍为负且没有充分披露的营运资本解释，则维持观察；只有实际H1和后续现金共同支持，才重新打开目标模型。 \\
002379 & 先等实际H1收入/OCF与预告归母；再取得调整后利润与归母的调节，并在后续季度检查铝价、氧化铝成本和分红。 & 监测规则：没有实际H1 OCF或口径调节即不发布范围；两项补齐后才重新评估。 \\
300014 & 先等实际H1收入、OCF和利润率；再在Q3检查储能/电池出货与现金，并调节CNY77/CNY103背后的EPS分母。 & 监测规则：H1/Q3现金未验证或分母无法调节即维持观察；不得取两个外部锚的平均数。 \\
000703 & 已有实际H1，下一步检查H2炼化利润、营运资本释放与OCF；同时观察价格是否重新形成可解释的安全边际。 & 监测规则：利润/现金任何一项转弱，或价格继续贴近一年高位，则不恢复低位资格。 \\
历史高空间样本 & 先以实际H1或后续季报补足缺口；再同时补足现金桥和当前可审计外部锚。 & 三项缺一不可；无当前目标或可复建锚时只允许监测。 \\
\bottomrule
\end{tabularx}

这些是研究监测规则，不是预测阈值或交易触发价。它们保证每次升级都有先后顺序：实际性在前，现金与口径在中，价格和外部锚在后。

\section{时间风险}

券商报告的评级通常对应其自身的6至12个月定义，并不自动等于2026年末目标。外部锚与实际财报的披露时点错位，也是本期不发布新目标价的重要原因。
